\section[The height equation and affine rigidity]{The renormalized height equation and affine rigidity}
\label{sec:fractional-rigidity}

The compactness lemma produces a single limiting height.  We now pass the
small Riesz pairing to that height, keeping the cancellation at infinity
throughout the passage.  The resulting equation is tested only against
functions of integral zero.  That class is sufficient to locate the Fourier
transform at the origin and hence to prove that the height is affine.
Interpreting a fractional operator modulo constants is natural for
functions of polynomial growth; compare
\cite{DipierroSavinValdinoci2019}. We prove the measurable weak
formulation needed here directly, rather than importing a regularity
or equivalence statement from a different formulation.

\subsection{A fixed normalization for the height pairing}

Write $k(z)=|z|^{-n-1}$ for $z\ne0$.  Let $\rho$ be a measure with upper
$n$-growth, and let $w$ be a measurable function with finite local
$L^2(\rho)$ norm and finite local energy
\[
 \iint_{B(0,T)^2}|w(x)-w(y)|^2 k(x-y)\,d\rho(x)d\rho(y)<\infty
 \qquad(T>0).
\]
Suppose also that
\begin{equation}\label{eq:height-weighted-tail}
 \int_{|y|>T}|w(y)|\,|y|^{-n-2}\,d\rho(y)<\infty
 \qquad(T>0).
\end{equation}
For a compactly supported Lipschitz function $\phi$, choose a bounded
measurable set $A$ containing a ball $B(0,R)$, with
$\supp\phi\subset B(0,H)$ and $2H<R$.  Define
\begin{align}
 \mathcal B_{\rho,A}(w,\phi)
 ={}&\frac12\iint_{A\times A}
       (w(x)-w(y))(\phi(x)-\phi(y))k(x-y)\,d\rho(x)d\rho(y)
       \label{eq:normalized-height-pairing}\\
 &+\iint_{A\times A^c}\phi(x)
       \bigl[w(x)k(x-y)-w(y)\{k(x-y)-k(-y)\}\bigr]
       \,d\rho(x)d\rho(y).\notag
\end{align}
This expression is defined even when $\int\phi\,d\rho\ne0$.  The set $A$
and the subtraction point $0$ are part of its normalization.

Both integrals in \eqref{eq:normalized-height-pairing} are absolutely
convergent.  For the interior integral, Cauchy--Schwarz, the Lipschitz bound,
and upper growth give, for every $h>0$,
\begin{align}
 &\iint_{\substack{x,y\in A\\|x-y|\le h}}
   |w(x)-w(y)|\,|\phi(x)-\phi(y)|k(x-y)\,d\rho(x)d\rho(y)
   \label{eq:height-diagonal-estimate}\\
 &\hspace{15mm}\le
  C\,\Lip(\phi)
  \left(\iint_{A^2}|w(x)-w(y)|^2k(x-y)\,d\rho(x)d\rho(y)\right)^{1/2}
  \rho(A)^{1/2}h^{1/2}.\notag
\end{align}
Indeed, the square of the remaining factor is bounded by
\[
 \Lip(\phi)^2\iint_{A^2\cap\{|x-y|\le h\}}|x-y|^{1-n}
 \,d\rho(x)d\rho(y)\le C\Lip(\phi)^2\rho(A)h.
\]
The part separated from the diagonal follows from local $L^1$ integrability.
For $x\in\supp\phi$ and $|y|\ge R$, the mean value theorem gives
\begin{equation}\label{eq:height-kernel-tail}
 k(x-y)\le C|y|^{-n-1},\qquad
 |k(x-y)-k(-y)|\le C H|y|^{-n-2}.
\end{equation}
The first term is integrable by upper growth and
$\phi w\in L^1(\rho)$; the second uses
\eqref{eq:height-weighted-tail}.

When $\int\phi\,d\rho=0$, the value in
\eqref{eq:normalized-height-pairing} is independent of the sufficiently
large localization $A$.  To see this without subtracting two divergent
integrals, restrict $\rho$ first to a common finite outer ball.  Symmetry
splits the half double integral into its $A^2$ part and the cross term.
The added center term integrates to
\[
 \left(\int\phi\,d\rho\right)
 \left(\int_{A^c}w(y)k(-y)\,d\rho(y)\right)=0
\]
in that finite restriction.  For two choices of $A$, both expressions
therefore equal the same finite double integral.  Now increase the outer
ball; \eqref{eq:height-kernel-tail} and
\eqref{eq:height-weighted-tail} permit dominated convergence of each
renormalized exterior term.  This proves the asserted independence.

There is also an exact connection with the vector pairing of
Lemma~\ref{lem:pairings}.  Suppose
\begin{equation}\label{eq:affine-source-height}
 u(x)-u(y)=e\cdot(x-y),\qquad w=u/\delta,\qquad \delta>0.
\end{equation}
Then for a compact Lipschitz test of $\rho$-mean zero,
\begin{equation}\label{eq:height-riesz-identification}
 \mathcal B_{\rho,A}(w,\phi)
       =\delta^{-1} e\cdot\mathcal R_\rho(\phi).
\end{equation}
The interior identity follows directly from
\eqref{eq:affine-source-height}.  In the exterior, subtraction of
$e\cdot K_n(-y)$ produces
\[
 \delta^{-1}\bigl[u(x)k(x-y)
       -u(y)\{k(x-y)-k(-y)\}-u(0)k(-y)\bigr].
\]
The last term integrates to zero against $\phi$.  Its exterior kernel is
integrable by upper growth and separation from $0$, so this cancellation is
legitimate.  The identity imposes no bound on $w(0)$.

\subsection{Passing to the same limiting height}

\begin{proposition}[The limiting height equation]\label{prop:height-equation}
Assume the hypotheses and conclusions of
Lemmas~\ref{lem:height-energy} and~\ref{lem:moving-compactness} for a scalar
height. We retain their notation and record the conditions used below.
Let $\mu_j$ converge on compact tests to
$\nu=q\H^n|_L$, where $q>0$ and $L$ is an $n$-plane through $0$.
Here $\H^n|_L$ has Euclidean normalization, equal to Lebesgue measure
in orthonormal coordinates on $L$.
Let $u_j$ satisfy \eqref{eq:affine-source-height} with $|e_j|\le1$, let
$w_j=u_j/\delta_j$, and let $v$ be the common limit given by
Lemma~\ref{lem:moving-compactness}.  Suppose $\delta_j>0$, $\delta_j\to0$,
$K_j\to\infty$,
and, for some fixed $0<\gamma<1$ and $M<\infty$,
\begin{equation}\label{eq:height-source-moment}
 \int_{B(0,2^k)}|u_j|^2\,d\mu_j
       \le M\delta_j^2\,2^{k(n+2+2\gamma)}
       \quad(0\le k\le K_j),\qquad
 \delta_j^{-1}2^{-K_j}\longrightarrow0.
\end{equation}
Assume that there are radii $A_j\to\infty$ such that, for every $j$ and
every compact Lipschitz test $\xi$ supported in $B(0,A_j)$ with
$\int\xi\,d\mu_j=0$,
\begin{equation}\label{eq:height-source-oscillation}
 |\mathcal R_{\mu_j}(\xi)|
      \le C(\Lip\xi+1)\delta_j^3.
\end{equation}
The constant here is independent of $j$ and of the test.

Then $v$ has finite local fractional energy and satisfies
\begin{equation}\label{eq:limit-height-weight}
 \int_L |v(x)|(1+|x|)^{-n-2}\,d\H^n(x)<\infty.
\end{equation}
After identifying $L$ isometrically with $\R^n$, for every compactly
supported smooth complex function $g$ with $\int_{\R^n}g=0$,
\begin{equation}\label{eq:limit-fractional-equation}
 \int_{\R^n}v(x)\mathcal Lg(x)\,dx=0,
 \qquad
 \mathcal Lg(x)=\frac12\int_{\R^n}
       \frac{2g(x)-g(x+h)-g(x-h)}{|h|^{n+1}}\,dh.
\end{equation}
The integral on the left is absolutely convergent.  The same conclusions
hold simultaneously for the finite family of normal heights in
Lemma~\ref{lem:moving-compactness}, without further extraction.
\end{proposition}

\begin{proof}
We retain the subsequence already selected in
Lemma~\ref{lem:moving-compactness}.  In the proof below every bounded
region and every test is chosen after this subsequence and its height $v$.

\emph{Local energy and singular pairings.}
Let $\Omega$ be one of the bounded continuity regions used in that lemma.
On its common compact closure, the three measures
\[
 \mu_j|_\Omega,\qquad w_j\mu_j|_\Omega,
       \qquad w_j^2\mu_j|_\Omega
\]
converge against continuous functions to the corresponding measures
$\nu|_\Omega$, $v\nu|_\Omega$, and $v^2\nu|_\Omega$.
Their total variations are uniformly bounded.  Finite sums of products of
continuous functions approximate a continuous function on the product of
this compact set uniformly.  Thus, after replacing $k(x-y)$ by
$(\max(h,|x-y|))^{-n-1}$, both the quadratic energy and every bilinear
pairing against a fixed bounded Lipschitz test converge.  For the energy,
expand $(w_j(x)-w_j(y))^2$; its three terms use exactly the three measures
above.  Monotone convergence as $h\downarrow0$ shows that the limiting
energy is no larger than the common bound from
Lemma~\ref{lem:height-energy}.  Positive-dimensional upper growth makes
the diagonal null.

Estimate~\eqref{eq:height-diagonal-estimate} applies uniformly to the
source and limiting heights on $\Omega$.  The error caused by capping the
bilinear kernel is bounded by its absolute integral on $|x-y|\le h$.
Consequently
\begin{align}\label{eq:height-local-bilinear-limit}
 &\iint_{\Omega^2}(w_j(x)-w_j(y))(\phi(x)-\phi(y))k(x-y)
       \,d\mu_j(x)d\mu_j(y)\\
 &\qquad\longrightarrow
 \iint_{\Omega^2}(v(x)-v(y))(\phi(x)-\phi(y))k(x-y)
       \,d\nu(x)d\nu(y).\notag
\end{align}
This argument first fixes $h$, then takes $j\to\infty$, and finally takes
$h\downarrow0$.

\emph{Inherited moments and the two exterior errors.}
The weak convergence of $w_j^2\mu_j$ on an enclosing continuity region,
tested with nonnegative cutoffs, passes
\eqref{eq:height-source-moment} to the limit at every fixed scale.  Changing
the radius by at most a factor two gives
\begin{equation}\label{eq:limiting-height-moments}
 \int_{L\cap B(0,t)}|v|^2\,d\nu\le C t^{n+2+2\gamma}
 \qquad(t\ge1).
\end{equation}
Each fixed radius eventually lies below $2^{K_j}$; no estimate past that
horizon for a fixed sample has been used.  Cauchy--Schwarz on dyadic shells,
followed by summation of $2^{k(\gamma-1)}$, now gives
\begin{equation}\label{eq:limiting-height-tail-bound}
 \int_{L\cap\{|y|>T\}} |v(y)|\,|y|^{-n-2}\,d\nu(y)
       \le C T^{\gamma-1}\qquad(T\ge2).
\end{equation}
Together with local $L^2$ integrability and $q>0$, this proves
\eqref{eq:limit-height-weight}.

Fix $H$ containing the supports of the tests to be considered.  Their
uniform supremum bounds and the local second-moment bound imply uniform
bounds on
$\int|\phi|\,d\mu_j$ and $\int|\phi w_j|\,d\mu_j$.
For $2H<T\le 2^{K_j}$, the same shell calculation using
\eqref{eq:height-source-moment} controls the renormalized exterior
integral outside $B(0,T)$ by
\begin{equation}\label{eq:two-height-tail-errors}
 C_\phi\bigl(T^{-1}+T^{\gamma-1}
                    +\delta_j^{-1}2^{-K_j}\bigr).
\end{equation}
Here is the bound beyond the last controlled radius.  The affine functions
$u_j$ are $1$-Lipschitz.  The mass of $B(0,1)$ is bounded below eventually,
by convergence to the plane measure, so
\eqref{eq:height-source-moment} implies
$|u_j(0)|\le1+C\delta_j$.  Hence $|u_j(y)|\le C+|y|$.
Upper growth gives
\[
 \int_{|y|\ge2^{K_j}}|w_j(y)||y|^{-n-2}\,d\mu_j(y)
      \le C\delta_j^{-1}2^{-K_j}.
\]
This bound uses the unnormalized height at $0$; the normalized value
$w_j(0)$ may diverge.  The $T^{-1}$ term in
\eqref{eq:two-height-tail-errors} comes from the $w_j(x)k(x-y)$ term,
whose $x$-integral stays on the fixed test support.  The kernel difference
in \eqref{eq:height-kernel-tail} gives the remaining terms.  Thus the order
is to choose $T$ large and then choose $j$ large.  The analogous limiting
tail tends to zero by \eqref{eq:limiting-height-tail-bound}.

\emph{A common normalization and mean correction.}
Take two fixed compact ambient Lipschitz functions $\phi,\eta$ such that
\[
 \int\phi\,d\nu=0,\qquad \int\eta\,d\nu\ne0.
\]
Choose one continuity region $A$ containing $B(0,R)$ with a positive
margin beyond both supports.  For each of these tests separately,
\begin{equation}\label{eq:fixed-normalization-convergence}
 \mathcal B_{\mu_j,A}(w_j,\phi)
       \longrightarrow\mathcal B_{\nu,A}(v,\phi),
 \qquad
 \mathcal B_{\mu_j,A}(w_j,\eta)
       \longrightarrow\mathcal B_{\nu,A}(v,\eta).
\end{equation}
For the interior terms this is
\eqref{eq:height-local-bilinear-limit}.  For the exterior terms first
restrict $y$ to a larger fixed continuity region $B$.  The kernels are
separated from their singularities on the nonzero test support.  To justify
the restrictions, let $F_h$ be the closed $h$-neighborhood of
$\partial A\cup\partial B$, for small $h>0$, and choose one bounded
continuity region containing all these collars.  The local second-moment
bound, Cauchy--Schwarz, and Portmanteau give
\[
 \limsup_{j\to\infty}\int_{F_h}|w_j|\,d\mu_j
 \le C\bigl(\limsup_{j\to\infty}\mu_j(F_h)\bigr)^{1/2}
 \le C\nu(F_h)^{1/2}\longrightarrow0
 \qquad(h\downarrow0).
\]
The final limit uses the $\nu$-nullity of both boundaries; the same
conclusion holds for ordinary mass.  Thus continuous cutoffs approximate
both restrictions, with product errors controlled by these collar masses
and the other factor's bounded total variation.  Convergence of $\mu_j$
and $w_j\mu_j$, followed by uniform approximation of the continuous
kernels by finite sums of product functions, now gives the exterior
limit on $B$.  Finally,
\eqref{eq:two-height-tail-errors} and
\eqref{eq:limiting-height-tail-bound} remove the outer restriction.
Both convergences in \eqref{eq:fixed-normalization-convergence} concern
the normalization \eqref{eq:normalized-height-pairing}; no Riesz pairing
is assigned to a test of nonzero mean.

Set
\[
 c_j=\frac{\int\phi\,d\mu_j}{\int\eta\,d\mu_j},
       \qquad \xi_j=\phi-c_j\eta.
\]
The denominator is nonzero eventually, and compact-test convergence gives
$c_j\to0$.  Eventually $|c_j|\le1$, so $\xi_j$ has a common compact
support, a uniform Lipschitz bound, and zero mean against the whole
$\mu_j$.  By linearity with the same $A$ and the same subtraction point,
\[
 \mathcal B_{\mu_j,A}(w_j,\xi_j)
       =\mathcal B_{\mu_j,A}(w_j,\phi)
        -c_j\mathcal B_{\mu_j,A}(w_j,\eta)
       \longrightarrow\mathcal B_{\nu,A}(v,\phi).
\]
The bump pairing on the right is already normalized and converges; the
only required fact about its coefficient is $c_j\to0$.
On the other hand, \eqref{eq:height-riesz-identification} and
\eqref{eq:height-source-oscillation} give
$|\mathcal B_{\mu_j,A}(w_j,\xi_j)|\le C_{\phi,\eta}\delta_j^2\to0$.
We conclude that
\begin{equation}\label{eq:compact-renormalized-height-equation}
 \mathcal B_{\nu,A}(v,\phi)=0
 \qquad\left(\int\phi\,d\nu=0\right).
\end{equation}

\emph{Intrinsic tests and the fractional operator.}
Let $U:\R^n\to L$ be a linear isometry.  For a compactly supported
intrinsic Lipschitz function $g$, choose an ambient radial cutoff $\chi$
equal to one on a ball containing $U(\supp g)$, and put
$\phi(x)=\chi(x)g(U^*x)$.
This is a compact ambient Lipschitz test and
$\phi(Ut)=g(t)$ for every $t$.  Its plane integral is $q\int g$.
A nonnegative compact radial bump positive at $0$ supplies the required
$\eta$.  Pulling \eqref{eq:compact-renormalized-height-equation} back by
$U$ and canceling the positive factor $q^2$ therefore gives the same
renormalized equation for $v(Ut)$ against every compact intrinsic
mean-zero test.  Localization independence allows us to take
$A=B(0,R)$ in the intrinsic plane.  We henceforth write this intrinsic
height simply as $v$.

For completeness, we identify this equation with
\eqref{eq:limit-fractional-equation}.  Let $g$ be compactly supported,
smooth and of integral zero, first real-valued, with support in $B(0,H)$;
choose $R>2H$.  For $\varepsilon>0$, define $\mathcal L_\varepsilon$
by restricting the $h$-integral in \eqref{eq:limit-fractional-equation}
to $|h|\ge\varepsilon$.  Reflection and translation give the absolutely
convergent single-increment formula
\[
 \mathcal L_\varepsilon g(x)
    =\int_{|x-y|\ge\varepsilon}(g(x)-g(y))k(x-y)\,dy.
\]
When $\varepsilon<R-H$, symmetrization on $B(0,R)^2$ yields
\begin{align*}
 \int v\mathcal L_\varepsilon g
 ={}&\frac12\iint_{B(0,R)^2\cap\{|x-y|\ge\varepsilon\}}
       (v(x)-v(y))(g(x)-g(y))k(x-y)\,dx\,dy\\
 &+\iint_{B(0,R)\times B(0,R)^c}g(x)
       [v(x)k(x-y)-v(y)\{k(x-y)-k(-y)\}]\,dx\,dy.
\end{align*}
To justify the exterior equality, integrate the test variable first:
$\int g=0$ permits the subtraction of $k(-y)$, after which
\eqref{eq:height-kernel-tail} makes the product integrable.  The direct
term involving $v(x)$ is separately integrable.  This proves the displayed
identity without applying Fubini to the uncancelled exterior height term.

The interior limit as $\varepsilon\downarrow0$ follows from
\eqref{eq:height-diagonal-estimate}.  On bounded $x$, the second difference
of $g$ is bounded by $C|h|^2$ near $h=0$, giving an
$\varepsilon$-independent bound for $\mathcal L_\varepsilon g$.
Outside a fixed ball containing the support, the cutoff is irrelevant and
$\int g=0$ gives
$|\mathcal L_\varepsilon g(x)|\le C_g(1+|x|)^{-n-2}$ for
$0<\varepsilon\le1$.  Dominated convergence using
\eqref{eq:limit-height-weight} therefore applies also to the left-hand
side.  The resulting identity is
$\int v\mathcal Lg=\mathcal B_{dx,B(0,R)}(v,g)=0$.
Real and imaginary parts give the complex statement.  All steps use the
same $v$ and the same preselected subsequence, so they apply to every
coordinate in a finite normal frame simultaneously.
\end{proof}

\subsection{Affine rigidity under the weighted tail bound}

We now prove the analytic rigidity statement needed for the limiting
heights.  Its tail hypothesis allows linear growth.  Write
$\mathcal S(\R^n)$ for the complex Schwartz space, with its usual
seminorms given by weighted suprema of derivatives.  On general Schwartz
tests, the same operator is given by
\begin{equation}\label{eq:rigidity-fractional-operator}
 \mathcal Lg(x)
 =-\frac12\int_{\R^n}
   \frac{g(x+h)+g(x-h)-2g(x)}{|h|^{n+1}}\,dh.
\end{equation}
The integral is absolutely convergent: the second difference is
$O(|h|^2)$ near zero and is bounded for large $h$.  The value assigned to
the integrand at $h=0$ has no effect.

\begin{lemma}[Weighted affine rigidity]\label{lem:fractional-rigidity}
Let $n\ge1$, and let $v:\R^n\to\R$ be measurable with
\begin{equation}\label{eq:rigidity-weighted-tail}
 \int_{\R^n}|v(x)|(1+|x|)^{-n-2}\,dx<\infty.
\end{equation}
Suppose that, for every $g\in C_c^\infty(\R^n;\mathbb C)$ with
$\int g=0$,
\begin{equation}\label{eq:rigidity-compact-equation}
 \int_{\R^n}v(x)\mathcal Lg(x)\,dx=0.
\end{equation}
Then there are $b\in\R$ and a real linear functional
$\ell:\R^n\to\R$ such that $v(x)=b+\ell(x)$ almost everywhere.
If $v$ is also essentially bounded, then $\ell=0$.
\end{lemma}

\begin{proof}
We first justify the extension of the test class in
\eqref{eq:rigidity-compact-equation}. We will then compute the positive
Fourier multiplier away from zero, deduce that the height is a
polynomial, and use the weighted tail bound to exclude degrees at
least two. For a Schwartz function $g$ put
\[
 p(g)=\sup_{y\in\R^n}(1+|y|)^{n+3}
       \bigl(|g(y)|+\|D^2g(y)\|\bigr),
 \qquad
 M_1(g)=\int_{\R^n}|y|\,|g(y)|\,dy.
\]
Here $D^2g$ is the second real derivative, with its operator norm.
The quantity $p$ is a continuous Schwartz seminorm, and
$M_1(g)\le C_n p(g)$, since
$\int |y|(1+|y|)^{-n-3}\,dy<\infty$.
We claim that
\begin{equation}\label{eq:rigidity-schwartz-decay}
 |\mathcal Lg(x)|\le C_n p(g)(1+|x|)^{-n-2}
 \qquad\text{if }\int g=0.
\end{equation}
Splitting the increment integral at radius one gives
$|\mathcal Lg(x)|\le C_n p(g)$ for every $x$.
For the spatial decay take $r=|x|\ge2$ and $R=r/2$.
If $|h|<R$, the segments from $x$ to $x\pm h$ stay outside
$B(0,r/2)$.  Taylor's formula therefore bounds the near part of
\eqref{eq:rigidity-fractional-operator} in absolute value by
\[
 C_n p(g)r^{-n-3}
      \int_{|h|<R}|h|^{1-n}\,dh
 \le C_n p(g)r^{-n-2}.
\]

For the far part let $E_x=\{y:|x-y|\ge R\}$ and
$k(z)=|z|^{-n-1}$.  Translation and reflection give exactly
\[
 g(x)\int_{|h|\ge R}k(h)\,dh
       -\int_{E_x}k(x-y)g(y)\,dy.
\]
These integrals are separately absolutely convergent.  The first term is
bounded by $C_n|g(x)|/r$.  Using the mean-zero condition in the second
term gives the identity
\begin{equation}\label{eq:rigidity-far-correction}
 \int_{E_x}k(x-y)g(y)\,dy
 =\int_{E_x}\bigl(k(x-y)-k(x)\bigr)g(y)\,dy
       -k(x)\int_{E_x^c}g(y)\,dy.
\end{equation}
On $|y|\le r/2$, the mean value theorem gives
\[
 |k(x-y)-k(x)|\le C_n|y|r^{-n-2}.
\]
On $E_x\cap\{|y|>r/2\}$, the kernel difference is bounded by
$C_n r^{-n-1}$, whereas
\[
 \int_{|y|>r/2}|g(y)|\,dy\le\frac{2}{r}M_1(g).
\]
Finally $E_x^c\subset\{|y|>r/2\}$, so the last term in
\eqref{eq:rigidity-far-correction} has the same bound
$C_n r^{-n-2}M_1(g)$.  This includes the correction for the omitted ball
around $x$.  The near and far estimates prove
\eqref{eq:rigidity-schwartz-decay}.

It follows from \eqref{eq:rigidity-weighted-tail} and
\eqref{eq:rigidity-schwartz-decay} that $v\mathcal Lg$ is integrable
for every mean-zero Schwartz function, and
\begin{equation}\label{eq:rigidity-pairing-continuity}
 \left|\int v\mathcal Lg\right|
 \le C_n p(g)\int |v(x)|(1+|x|)^{-n-2}\,dx.
\end{equation}
Choose $\chi,\rho\in C_c^\infty(\R^n)$ with $\chi=1$ near zero and
$\int\rho=1$.  For a mean-zero Schwartz function $g$ set
\[
 g_j(x)=\chi(x/j)g(x)
        -\left(\int\chi(y/j)g(y)\,dy\right)\rho(x).
\]
Each $g_j$ is compactly supported and has mean zero.
The product rule and rapid decrease show
$\chi(\cdot/j)g\to g$ in every Schwartz seminorm; also
$\int\chi(y/j)g(y)\,dy\to\int g=0$.
Thus $g_j\to g$ in Schwartz space.
Applying \eqref{eq:rigidity-pairing-continuity} to $g_j-g$ extends
\eqref{eq:rigidity-compact-equation} to every mean-zero
$g\in\mathcal S(\R^n)$.

We next compute the Fourier multiplier needed for localization away from
zero.  Use the convention
\[
 \widehat g(\xi)=\int_{\R^n}e^{-2\pi i x\cdot\xi}g(x)\,dx,
 \qquad
 \mathcal F^{-1}g(x)=\int_{\R^n}e^{2\pi i x\cdot\xi}g(\xi)\,d\xi.
\]
Define
\[
 a_n(\xi)=\int_{\R^n}
        \frac{1-\cos(2\pi\xi\cdot h)}{|h|^{n+1}}\,dh.
\]
This is an absolutely convergent nonnegative integral: quadratic
cancellation controls the origin, and boundedness of $1-\cos$ controls
infinity.  If $\xi\ne0$, the numerator equals $2$ at
$h=\xi/(2|\xi|^2)$ and is positive on a neighborhood of that point.
Consequently $a_n(\xi)>0$.  Orthogonal changes of variables show that
$a_n$ is radial, and the substitution $u=th$ gives
$a_n(t\xi)=t a_n(\xi)$ for $t>0$.  Hence
\begin{equation}\label{eq:rigidity-positive-symbol}
 a_n(\xi)=c_n|\xi|,\qquad c_n=a_n(e_1)>0.
\end{equation}

For $g\in\mathcal S(\R^n)$ the full absolute integral required to
interchange frequency and increments is finite:
\[
 \int_{\R^n}\int_{\R^n}
   \frac{1-\cos(2\pi\xi\cdot h)}{|h|^{n+1}}|g(\xi)|\,dh\,d\xi
 =c_n\int_{\R^n}|\xi|\,|g(\xi)|\,d\xi<\infty.
\]
Fubini's theorem and the second difference of the exponential therefore
give, pointwise in $x$,
\begin{equation}\label{eq:rigidity-fourier-multiplier}
 \mathcal L(\mathcal F^{-1}g)(x)
   =c_n\int_{\R^n}e^{2\pi i x\cdot\xi}|\xi|g(\xi)\,d\xi.
\end{equation}
The factor $-2$ in that second difference cancels the factor $-1/2$ in
\eqref{eq:rigidity-fractional-operator}.

The weighted hypothesis implies $v\in L^1_{\mathrm{loc}}$ and defines a
tempered distribution
\[
 T_v(g)=\int_{\R^n}v(x)g(x)\,dx.
\]
Indeed this integral is bounded by
$\sup_x(1+|x|)^{n+2}|g(x)|$ times the integral in
\eqref{eq:rigidity-weighted-tail}.  All distributional pairings here are
complex bilinear, without complex conjugation.
Let $S=\mathcal F^{-1}T_v$, so
$S(g)=T_v(\mathcal F^{-1}g)$.
Take $\psi\in C_c^\infty(\R^n\setminus\{0\})$.
The function $q(\xi)=\psi(\xi)/|\xi|$, extended by zero near the
origin, is smooth and compactly supported.  Its inverse Fourier transform
is Schwartz and has integral $q(0)=0$.  The extended weak equation and
\eqref{eq:rigidity-fourier-multiplier} give
\[
 0=T_v\bigl(\mathcal L(\mathcal F^{-1}q)\bigr)
   =c_nT_v(\mathcal F^{-1}\psi)=c_n S(\psi).
\]
Thus $S$ is supported at the origin.

For completeness we prove the point-support consequence; compare
\cite[Theorem~2.3.4]{Hormander2003}. Continuity of
$S$ on Schwartz space supplies $N\in\mathbb N$ and $C<\infty$ such that
\begin{equation}\label{eq:rigidity-finite-order}
 |S(g)|\le C\sum_{|\alpha|\le N}
       \sup_x(1+|x|)^N|\partial^\alpha g(x)|.
\end{equation}
A Schwartz function vanishing near zero is annihilated by $S$: multiply it
by expanding compact cutoffs, use the support property, and pass to the
Schwartz limit.  Consequently $S$ has the same value on two Schwartz
functions agreeing near zero.

Suppose now that $\partial^\alpha g(0)=0$ for every $|\alpha|\le N$.
Choose $\eta\in C_c^\infty(B(0,2))$ equal to one near zero and put
$\eta_\varepsilon(x)=\eta(x/\varepsilon)$.
For $0<\varepsilon<1/2$, Taylor's formula and the product rule give
\[
 \sup_x(1+|x|)^N
       |\partial^\alpha(\eta_\varepsilon g)(x)|
       \le C_g\varepsilon^{N+1-|\alpha|}
       \le C_g\varepsilon
       \qquad (|\alpha|\le N).
\]
To see the power, a derivative of order $j$ falling on the cutoff costs
$\varepsilon^{-j}$, while the remaining derivative of $g$, of order
$|\alpha|-j$, is $O(\varepsilon^{N+1-|\alpha|+j})$ on its support.
Since $S(g)=S(\eta_\varepsilon g)$, letting $\varepsilon\downarrow0$
in \eqref{eq:rigidity-finite-order} proves $S(g)=0$.

It follows that $S$ depends only on the derivatives through order $N$ at
zero.  More explicitly, subtract from $g$ the function
\[
 \eta(x)\sum_{|\alpha|\le N}
           \frac{\partial^\alpha g(0)}{\alpha!}x^\alpha.
\]
The remainder has vanishing jet through order $N$ at zero, and therefore there are
constants $c_\alpha\in\mathbb C$ such that
\[
 S(g)=\sum_{|\alpha|\le N}c_\alpha\partial^\alpha g(0).
\]
Using $T_v(g)=S(\widehat g)$ and differentiating the Fourier integral
of a Schwartz function at zero gives
\[
 T_v(g)=\int_{\R^n}P(x)g(x)\,dx,
 \qquad
 P(x)=\sum_{|\alpha|\le N}c_\alpha(-2\pi i x)^\alpha.
\]
Every polynomially weighted Schwartz function is integrable, so the
finite sum and these integrals may be interchanged.
Since $v$ is locally integrable, equality on compact smooth tests implies
$v=P$ almost everywhere.

Finally, \eqref{eq:rigidity-weighted-tail} implies
\[
 \int_{|x|\ge1}|P(x)|\,|x|^{-n-2}\,dx<\infty.
\]
If $P$ had degree $k\ge2$, its nonzero leading homogeneous part would
have absolute value bounded below on a nonempty open subset of the unit
sphere.  On the corresponding cone and for all sufficiently large $r$,
$|P(r\theta)|\ge c r^k$.  The last integral would then dominate a
positive constant times
\[
 \int_{R_0}^{\infty}r^{k-3}\,dr,
\]
which diverges for $k\ge2$.  (For $n=1$, one uses either of the two
rays.)  Hence $P$ has degree at most one.  Taking real parts gives
$v=b+\ell$ almost everywhere.
If $|v|\le M$ almost everywhere, continuity gives
$|b+\ell(x)|\le M$ for every $x$, since any open set violating this bound
would have positive measure.  Evaluating along rays shows $\ell=0$.
\end{proof}

The same conclusion holds if the compact weak equation is initially
known only for real mean-zero tests: apply it separately to the real and
imaginary parts, using the absolute integrability proved above.
